\documentclass[../../Main.tex]{subfiles}
\begin{document}
\section{自定义GUI}
在模板中，GUI的定义被分为了三部分，一部分是GUI的颜色，定义在gui.rpy文件中；一部分则是GUI的布局，通过界面语言定义在screen.rpy文件中（在下一章我们将会学习界面语言）；还有一部分是图片等已经生成的资源，存储在images.rpa文件中。其中对颜色的修改主要涉及gui.rpy文件与images.rpa文件中。

在gui.rpy中，主要定义了以下几个方面的内容：
\begin{itemize}
    \item GUI配置变量：包括字体、音效等。
    \item 全局颜色（这里的颜色仅包括：文字、强调色，对于GUI界面背景等需要进行解包处理）；
    \item 文本大小、选项宽窄等。
\end{itemize}

对于gui.rpy文件的修改是比较简单的，直接打开文件进行修改即可。需要注意的是，修改颜色时应当使用十六进制颜色码，修改字体时您需要参照第\ref{4}章的内容将字体文件添加在mod\_assets/文件夹中并进行定义引用。

\begin{Warning}
    我们不推荐您修改分辨率与文本框等内容（除非您需要将您的模组移植到手机端等平台上），因为这可能会导致一些意想不到的问题。
\end{Warning}

对于gui.rpy文件的修改，您可以参考\url{https://doc.renpy.cn/zh-CN/gui.html}中的内容。

\subparagraph{解包rpa}
对于拥有Python环境的读者，您可以使用unrpa工具。对于大部分读者，您可以通过rpaextract工具进行解包。
您可以在\url{https://iwanplays.itch.io/rpaex}下载rpaextract工具。rpaextract工具使用非常简单，下载后将您需要解包的rpa文件拖入rpaextract.exe程序中即可进行解包。

解包images.rpa文件后，您可以在文件夹中找到gui文件夹，里面存储了大量的GUI资源图片。您可以根据需要使用软件对其进行修改与替换。

请注意，所有修改后的GUI资源图片均需要放置在mod\_assets/gui/文件夹中，按照第\ref{4}章与\url{https://doc.renpy.cn/zh-CN/gui.html}的内容进行定义引用。

\end{document}